% =============================================================================
% Informe técnico — Sistema Master Gateway
% Uso:
%   1) Compilar solo: pdflatex docs/informe_proyecto.tex
%   2) Incorporar en otro documento:
%        % Cuerpo del informe (sin preámbulo) para incorporar en otro documento LaTeX:
%   % Cuerpo del informe (sin preámbulo) para incorporar en otro documento LaTeX:
%   % Cuerpo del informe (sin preámbulo) para incorporar en otro documento LaTeX:
%   \input{docs/informe_proyecto_cuerpo}
% Requiere en el documento padre (según uso): babel spanish, booktabs, listings,
% enumitem, hyperref (opcionales pero recomendados).

\section{Introducción}
El proyecto académico \emph{Sistema de Autenticación y Autorización Centralizado
(Master Gateway)} responde a la necesidad de una puerta maestra de seguridad
para arquitecturas de microservicios. En lugar de que cada servicio gestione
usuarios y credenciales, el Master autentica usuarios, obliga a seleccionar un
rol activo, emite JWT con privilegio mínimo, expone validación interna Zero Trust
para microservicios hijos y administra usuarios, roles, módulos, menús y permisos
con soft delete y auditoría.

Principios rectores: Zero Trust, Least Privilege, Shift-Left Security,
autenticación stateless y soft delete.

\section{Objetivos}
\begin{enumerate}
  \item Implementar el flujo login $\rightarrow$ TempToken $\rightarrow$ selección de rol
        $\rightarrow$ AccessToken + RefreshToken rotativo.
  \item Modelar entidades con auditoría y estado ACTIVO/INACTIVO.
  \item Generar menú dinámico con CTE recursiva filtrado por rol JWT.
  \item Integrar un microservicio hijo (Ventas) bajo Zero Trust.
  \item Entregar SPA con rutas protegidas y navegación dinámica.
  \item Automatizar calidad y despliegue (GitHub Actions, SonarCloud, SAST/ML, Railway, Telegram).
\end{enumerate}

\section{Arquitectura del sistema}
\subsection{Vista lógica}
\begin{enumerate}
  \item \textbf{Master Gateway}: API Node.js/TypeScript (auth, RBAC, CRUD, validate-token).
  \item \textbf{PostgreSQL} vía Prisma ORM.
  \item \textbf{SPA React + Vite}: login, selector de rol, administración y Ventas.
  \item \textbf{Microservicio Ventas}: sin DB de usuarios; valida tokens en el Master.
\end{enumerate}

\subsection{Flujo Zero Trust}
El frontend envía el AccessToken a Ventas; Ventas llama a
\texttt{POST /api/internals/validate-token} con \texttt{x-internal-api-key}.
Solo si el Master confirma validez y permisos se ejecuta la lógica de negocio.

\subsection{Stack}
Master: Express, TypeScript, Prisma, Argon2, Jose (JWT), Helmet, Zod.
Frontend: React 19, Vite, React Router.
CI/CD: GitHub Actions, SonarCloud, SAST/ML, Railway CLI, Telegram.

\section{Modelo de datos}
Entidades: User, Role, UserRole, Module, RoleModule, Menu (Adjacency List),
RoleMenu, Permission, RolePermission, RefreshToken, TokenRevocation, AuditLog.
Menú: solo hojas con URL, anticiclos, exclusión de inactivos.

\section{Autenticación y autorización}
Login emite TempToken + roles; la SPA fuerza select-role; el AccessToken incluye
únicamente el rol seleccionado y sus permisos. Refresh rotativo y logout con
revocación por \texttt{jti}. Endpoint interno
\texttt{/api/internals/validate-token} responde
\texttt{active}, \texttt{userId}, \texttt{roleId}, \texttt{roleName}, \texttt{permissions}.

\section{Frontend}
SPA con sessionStorage, interceptor Bearer/refresh, menú dinámico, administración
de usuarios/roles, pantallas 403/sesión/token expirado y cancelación del selector
de rol para cambiar de usuario. La UI no decide autorización final.

\section{Microservicio Ventas}
Endpoints con permisos \texttt{VENTAS\_READ}/\texttt{VENTAS\_CREATE}; retries ante
cold start PaaS. Roles demo: ADMIN, VENDEDOR, AUDITOR.

\section{Seguridad}
Argon2, rate limiting, Helmet, Zod, soft delete, secretos en entorno, mensajes
de login genéricos. JWT HS256 + validate-token (RS256 opcional no bloqueante).

\section{DevSecOps}
Ramas \texttt{feature $\rightarrow$ dev $\rightarrow$ test $\rightarrow$ main}.
Pipeline en merge a main: build/tests, Sonar Quality Gate obligatorio, SAST/ML,
deploy Railway, Telegram. Despliegue recomendado: Railway (alternativa Render).

\section{Pruebas}
Vitest, smoke E2E (\texttt{npm run smoke}), checklist en
\texttt{docs/COMPLIANCE\_CHECKLIST.md}.

\section{Conclusiones}
Se cumple el alcance académico Full-Stack seguro (Master + SPA + Ventas + DevSecOps).
Pendiente operativo: secrets GitHub, branch protection, SonarCloud y PaaS.
Mejora opcional: JWT asimétrico.

% Requiere en el documento padre (según uso): babel spanish, booktabs, listings,
% enumitem, hyperref (opcionales pero recomendados).

\section{Introducción}
El proyecto académico \emph{Sistema de Autenticación y Autorización Centralizado
(Master Gateway)} responde a la necesidad de una puerta maestra de seguridad
para arquitecturas de microservicios. En lugar de que cada servicio gestione
usuarios y credenciales, el Master autentica usuarios, obliga a seleccionar un
rol activo, emite JWT con privilegio mínimo, expone validación interna Zero Trust
para microservicios hijos y administra usuarios, roles, módulos, menús y permisos
con soft delete y auditoría.

Principios rectores: Zero Trust, Least Privilege, Shift-Left Security,
autenticación stateless y soft delete.

\section{Objetivos}
\begin{enumerate}
  \item Implementar el flujo login $\rightarrow$ TempToken $\rightarrow$ selección de rol
        $\rightarrow$ AccessToken + RefreshToken rotativo.
  \item Modelar entidades con auditoría y estado ACTIVO/INACTIVO.
  \item Generar menú dinámico con CTE recursiva filtrado por rol JWT.
  \item Integrar un microservicio hijo (Ventas) bajo Zero Trust.
  \item Entregar SPA con rutas protegidas y navegación dinámica.
  \item Automatizar calidad y despliegue (GitHub Actions, SonarCloud, SAST/ML, Railway, Telegram).
\end{enumerate}

\section{Arquitectura del sistema}
\subsection{Vista lógica}
\begin{enumerate}
  \item \textbf{Master Gateway}: API Node.js/TypeScript (auth, RBAC, CRUD, validate-token).
  \item \textbf{PostgreSQL} vía Prisma ORM.
  \item \textbf{SPA React + Vite}: login, selector de rol, administración y Ventas.
  \item \textbf{Microservicio Ventas}: sin DB de usuarios; valida tokens en el Master.
\end{enumerate}

\subsection{Flujo Zero Trust}
El frontend envía el AccessToken a Ventas; Ventas llama a
\texttt{POST /api/internals/validate-token} con \texttt{x-internal-api-key}.
Solo si el Master confirma validez y permisos se ejecuta la lógica de negocio.

\subsection{Stack}
Master: Express, TypeScript, Prisma, Argon2, Jose (JWT), Helmet, Zod.
Frontend: React 19, Vite, React Router.
CI/CD: GitHub Actions, SonarCloud, SAST/ML, Railway CLI, Telegram.

\section{Modelo de datos}
Entidades: User, Role, UserRole, Module, RoleModule, Menu (Adjacency List),
RoleMenu, Permission, RolePermission, RefreshToken, TokenRevocation, AuditLog.
Menú: solo hojas con URL, anticiclos, exclusión de inactivos.

\section{Autenticación y autorización}
Login emite TempToken + roles; la SPA fuerza select-role; el AccessToken incluye
únicamente el rol seleccionado y sus permisos. Refresh rotativo y logout con
revocación por \texttt{jti}. Endpoint interno
\texttt{/api/internals/validate-token} responde
\texttt{active}, \texttt{userId}, \texttt{roleId}, \texttt{roleName}, \texttt{permissions}.

\section{Frontend}
SPA con sessionStorage, interceptor Bearer/refresh, menú dinámico, administración
de usuarios/roles, pantallas 403/sesión/token expirado y cancelación del selector
de rol para cambiar de usuario. La UI no decide autorización final.

\section{Microservicio Ventas}
Endpoints con permisos \texttt{VENTAS\_READ}/\texttt{VENTAS\_CREATE}; retries ante
cold start PaaS. Roles demo: ADMIN, VENDEDOR, AUDITOR.

\section{Seguridad}
Argon2, rate limiting, Helmet, Zod, soft delete, secretos en entorno, mensajes
de login genéricos. JWT HS256 + validate-token (RS256 opcional no bloqueante).

\section{DevSecOps}
Ramas \texttt{feature $\rightarrow$ dev $\rightarrow$ test $\rightarrow$ main}.
Pipeline en merge a main: build/tests, Sonar Quality Gate obligatorio, SAST/ML,
deploy Railway, Telegram. Despliegue recomendado: Railway (alternativa Render).

\section{Pruebas}
Vitest, smoke E2E (\texttt{npm run smoke}), checklist en
\texttt{docs/COMPLIANCE\_CHECKLIST.md}.

\section{Conclusiones}
Se cumple el alcance académico Full-Stack seguro (Master + SPA + Ventas + DevSecOps).
Pendiente operativo: secrets GitHub, branch protection, SonarCloud y PaaS.
Mejora opcional: JWT asimétrico.

% Requiere en el documento padre (según uso): babel spanish, booktabs, listings,
% enumitem, hyperref (opcionales pero recomendados).

\section{Introducción}
El proyecto académico \emph{Sistema de Autenticación y Autorización Centralizado
(Master Gateway)} responde a la necesidad de una puerta maestra de seguridad
para arquitecturas de microservicios. En lugar de que cada servicio gestione
usuarios y credenciales, el Master autentica usuarios, obliga a seleccionar un
rol activo, emite JWT con privilegio mínimo, expone validación interna Zero Trust
para microservicios hijos y administra usuarios, roles, módulos, menús y permisos
con soft delete y auditoría.

Principios rectores: Zero Trust, Least Privilege, Shift-Left Security,
autenticación stateless y soft delete.

\section{Objetivos}
\begin{enumerate}
  \item Implementar el flujo login $\rightarrow$ TempToken $\rightarrow$ selección de rol
        $\rightarrow$ AccessToken + RefreshToken rotativo.
  \item Modelar entidades con auditoría y estado ACTIVO/INACTIVO.
  \item Generar menú dinámico con CTE recursiva filtrado por rol JWT.
  \item Integrar un microservicio hijo (Ventas) bajo Zero Trust.
  \item Entregar SPA con rutas protegidas y navegación dinámica.
  \item Automatizar calidad y despliegue (GitHub Actions, SonarCloud, SAST/ML, Railway, Telegram).
\end{enumerate}

\section{Arquitectura del sistema}
\subsection{Vista lógica}
\begin{enumerate}
  \item \textbf{Master Gateway}: API Node.js/TypeScript (auth, RBAC, CRUD, validate-token).
  \item \textbf{PostgreSQL} vía Prisma ORM.
  \item \textbf{SPA React + Vite}: login, selector de rol, administración y Ventas.
  \item \textbf{Microservicio Ventas}: sin DB de usuarios; valida tokens en el Master.
\end{enumerate}

\subsection{Flujo Zero Trust}
El frontend envía el AccessToken a Ventas; Ventas llama a
\texttt{POST /api/internals/validate-token} con \texttt{x-internal-api-key}.
Solo si el Master confirma validez y permisos se ejecuta la lógica de negocio.

\subsection{Stack}
Master: Express, TypeScript, Prisma, Argon2, Jose (JWT), Helmet, Zod.
Frontend: React 19, Vite, React Router.
CI/CD: GitHub Actions, SonarCloud, SAST/ML, Railway CLI, Telegram.

\section{Modelo de datos}
Entidades: User, Role, UserRole, Module, RoleModule, Menu (Adjacency List),
RoleMenu, Permission, RolePermission, RefreshToken, TokenRevocation, AuditLog.
Menú: solo hojas con URL, anticiclos, exclusión de inactivos.

\section{Autenticación y autorización}
Login emite TempToken + roles; la SPA fuerza select-role; el AccessToken incluye
únicamente el rol seleccionado y sus permisos. Refresh rotativo y logout con
revocación por \texttt{jti}. Endpoint interno
\texttt{/api/internals/validate-token} responde
\texttt{active}, \texttt{userId}, \texttt{roleId}, \texttt{roleName}, \texttt{permissions}.

\section{Frontend}
SPA con sessionStorage, interceptor Bearer/refresh, menú dinámico, administración
de usuarios/roles, pantallas 403/sesión/token expirado y cancelación del selector
de rol para cambiar de usuario. La UI no decide autorización final.

\section{Microservicio Ventas}
Endpoints con permisos \texttt{VENTAS\_READ}/\texttt{VENTAS\_CREATE}; retries ante
cold start PaaS. Roles demo: ADMIN, VENDEDOR, AUDITOR.

\section{Seguridad}
Argon2, rate limiting, Helmet, Zod, soft delete, secretos en entorno, mensajes
de login genéricos. JWT HS256 + validate-token (RS256 opcional no bloqueante).

\section{DevSecOps}
Ramas \texttt{feature $\rightarrow$ dev $\rightarrow$ test $\rightarrow$ main}.
Pipeline en merge a main: build/tests, Sonar Quality Gate obligatorio, SAST/ML,
deploy Railway, Telegram. Despliegue recomendado: Railway (alternativa Render).

\section{Pruebas}
Vitest, smoke E2E (\texttt{npm run smoke}), checklist en
\texttt{docs/COMPLIANCE\_CHECKLIST.md}.

\section{Conclusiones}
Se cumple el alcance académico Full-Stack seguro (Master + SPA + Ventas + DevSecOps).
Pendiente operativo: secrets GitHub, branch protection, SonarCloud y PaaS.
Mejora opcional: JWT asimétrico.
   % (si se separa el cuerpo)
%      o copiar las secciones desde \section{Introducción} en adelante.
% =============================================================================
\documentclass[12pt,a4paper]{article}

\usepackage[utf8]{inputenc}
\usepackage[T1]{fontenc}
\usepackage[spanish]{babel}
\usepackage{geometry}
\usepackage{hyperref}
\usepackage{booktabs}
\usepackage{longtable}
\usepackage{array}
\usepackage{enumitem}
\usepackage{listings}
\usepackage{xcolor}
\usepackage{graphicx}
\usepackage{titlesec}
\usepackage{fancyhdr}
\usepackage{setspace}

\geometry{margin=2.5cm}
\onehalfspacing

\hypersetup{
  colorlinks=true,
  linkcolor=teal!60!black,
  urlcolor=teal!60!black,
  citecolor=teal!60!black
}

\definecolor{codebg}{RGB}{245,247,248}
\lstset{
  backgroundcolor=\color{codebg},
  basicstyle=\ttfamily\footnotesize,
  breaklines=true,
  frame=single,
  rulecolor=\color{black!20},
  columns=fullflexible
}

\pagestyle{fancy}
\fancyhf{}
\fancyhead[L]{\small Master Gateway --- Informe técnico}
\fancyhead[R]{\small Desarrollo de Software Seguro}
\fancyfoot[C]{\thepage}

\title{%
  Informe técnico del proyecto\\
  \large Sistema de Autenticación y Autorización Centralizado\\
  (Master Gateway)
}
\author{Equipo de desarrollo \\ Desarrollo de Software Seguro --- Parcial III}
\date{\today}

\begin{document}
\maketitle
\begin{abstract}
Este documento describe la arquitectura, el diseño de seguridad, los componentes
implementados y la estrategia DevSecOps del \textbf{Master Gateway}: un
microservicio maestro que centraliza autenticación, autorización basada en roles
(RBAC), menú dinámico y validación Zero Trust para microservicios hijos.
El sistema incluye una SPA React, un microservicio de Ventas y un pipeline
CI/CD con SonarCloud, análisis SAST/ML, notificaciones Telegram y despliegue
automático en Railway.
\end{abstract}

\tableofcontents
\newpage

% -----------------------------------------------------------------------------
\section{Introducción}
% -----------------------------------------------------------------------------
El proyecto académico \emph{Sistema de Autenticación y Autorización Centralizado
(Master Gateway)} responde a la necesidad de una puerta maestra de seguridad
para arquitecturas de microservicios. En lugar de que cada servicio gestione
usuarios y credenciales, el Master:

\begin{itemize}[leftmargin=1.5em]
  \item autentica usuarios;
  \item obliga a seleccionar un rol activo (workspace);
  \item emite JWT con \textbf{privilegio mínimo} (solo el rol elegido y sus permisos);
  \item expone validación interna para que los microservicios hijos verifiquen
        tokens sin confiar en el frontend;
  \item administra usuarios, roles, módulos, menús y permisos con soft delete y auditoría.
\end{itemize}

Principios rectores: \textbf{Zero Trust}, \textbf{Least Privilege},
\textbf{Shift-Left Security}, autenticación \textbf{stateless} y
\textbf{soft delete}.

% -----------------------------------------------------------------------------
\section{Objetivos}
% -----------------------------------------------------------------------------
\begin{enumerate}[leftmargin=1.5em]
  \item Implementar el flujo login $\rightarrow$ TempToken $\rightarrow$ selección de rol
        $\rightarrow$ AccessToken + RefreshToken rotativo.
  \item Modelar entidades con auditoría y estado \texttt{ACTIVO}/\texttt{INACTIVO}.
  \item Generar menú dinámico con CTE recursiva filtrado por rol JWT.
  \item Integrar al menos un microservicio hijo (Ventas) bajo Zero Trust.
  \item Entregar SPA con rutas protegidas y navegación dinámica.
  \item Automatizar calidad y despliegue (GitHub Actions, SonarCloud, SAST/ML, Railway, Telegram).
\end{enumerate}

% -----------------------------------------------------------------------------
\section{Arquitectura del sistema}
% -----------------------------------------------------------------------------
\subsection{Vista lógica}
El sistema se compone de cuatro bloques principales:

\begin{enumerate}[leftmargin=1.5em]
  \item \textbf{Master Gateway (API Node.js/TypeScript)}: autenticación,
        autorización, CRUD administrativo y endpoint interno de validación.
  \item \textbf{PostgreSQL}: persistencia vía Prisma ORM.
  \item \textbf{SPA (React + Vite)}: experiencia de usuario, selector de rol y
        administración.
  \item \textbf{Microservicio Ventas}: lógica de negocio sin base de usuarios;
        valida cada petición contra el Master.
\end{enumerate}

\subsection{Flujo de confianza (Zero Trust)}
\begin{enumerate}[leftmargin=1.5em]
  \item El usuario se autentica en el Master mediante la SPA.
  \item Tras seleccionar rol, la SPA obtiene un AccessToken.
  \item Al invocar Ventas, la SPA envía \texttt{Authorization: Bearer <token>}.
  \item Ventas \textbf{no confía} en el frontend: llama a
        \texttt{POST /api/internals/validate-token} con
        \texttt{x-internal-api-key} y el token del usuario.
  \item Solo si el Master confirma validez y permisos, Ventas ejecuta la operación.
\end{enumerate}

\subsection{Stack tecnológico}
\begin{center}
\begin{tabular}{@{}ll@{}}
\toprule
\textbf{Capa} & \textbf{Tecnología} \\
\midrule
Backend Master & Node.js, Express, TypeScript, Prisma, PostgreSQL \\
Seguridad & Argon2, Jose (JWT HS256), Helmet, rate limit, Zod \\
Frontend & React 19, Vite, React Router \\
Microservicio & Express + validación vía Master \\
CI/CD & GitHub Actions, SonarCloud, Python/ML SAST, Railway CLI, Telegram \\
\bottomrule
\end{tabular}
\end{center}

% -----------------------------------------------------------------------------
\section{Modelo de datos}
% -----------------------------------------------------------------------------
Entidades principales (todas con auditoría: \texttt{fechaCreacion},
\texttt{fechaActualizacion}, \texttt{creadoPor}, \texttt{actualizadoPor} y
\texttt{estado}):

\begin{itemize}[leftmargin=1.5em]
  \item \textbf{User}, \textbf{Role}, \textbf{UserRole}
  \item \textbf{Module}, \textbf{RoleModule}
  \item \textbf{Menu} (Adjacency List: \texttt{parentId}, \texttt{orden}, \texttt{icono}, \texttt{url})
  \item \textbf{RoleMenu}, \textbf{Permission}, \textbf{RolePermission}
  \item \textbf{RefreshToken} (rotación / \texttt{reemplazadoPor})
  \item \textbf{TokenRevocation} (logout e invalidación por \texttt{jti})
  \item \textbf{AuditLog} (login, rol, refresh, logout, usuarios, permisos, intentos no autorizados)
\end{itemize}

Reglas de menú:
\begin{itemize}[leftmargin=1.5em]
  \item solo nodos hoja pueden tener URL;
  \item se impiden ciclos en \texttt{parentId};
  \item nodos inactivos (y módulos inactivos) se excluyen del árbol.
\end{itemize}

% -----------------------------------------------------------------------------
\section{Autenticación y autorización}
% -----------------------------------------------------------------------------
\subsection{Flujo de autenticación}
\begin{enumerate}[leftmargin=1.5em]
  \item \texttt{POST /api/auth/login}: valida email/password (Argon2).
        Respuesta: \texttt{tempToken} de corta duración + roles activos.
  \item SPA muestra selector obligatorio de rol (no hay acceso directo al dashboard).
  \item \texttt{POST /api/auth/select-role}: valida TempToken, pertenencia del rol
        y estado activo. Emite:
        \begin{itemize}
          \item AccessToken JWT (vida corta);
          \item RefreshToken rotativo;
          \item rol activo y permisos.
        \end{itemize}
  \item \texttt{POST /api/auth/refresh-token}: rota el refresh; reutilización sospechosa
        implica revocación.
  \item \texttt{POST /api/auth/logout}: revoca refresh y registra \texttt{jti} del access.
\end{enumerate}

\subsection{Claims del AccessToken (Least Privilege)}
El JWT definitivo \textbf{no} incluye todos los roles del usuario. Contiene, entre otros:
\texttt{sub}, \texttt{roleId}, \texttt{roleName}, \texttt{permissions},
\texttt{iss}, \texttt{aud}, \texttt{iat}, \texttt{exp}, \texttt{jti}, \texttt{type=access}.

\subsection{Validación interna}
\texttt{POST /api/internals/validate-token} (protegido por API key interna) responde:
\begin{lstlisting}
{
  "active": true,
  "userId": "...",
  "roleId": "...",
  "roleName": "VENDEDOR",
  "permissions": ["VENTAS_READ", "VENTAS_CREATE"]
}
\end{lstlisting}
Opcionalmente exige \texttt{requiredPermissions}.

% -----------------------------------------------------------------------------
\section{API del Master (resumen)}
% -----------------------------------------------------------------------------
\begin{center}
\begin{tabular}{@{}p{6.2cm}p{7.5cm}@{}}
\toprule
\textbf{Grupo} & \textbf{Endpoints} \\
\midrule
Auth &
\texttt{/api/auth/login}, \texttt{select-role}, \texttt{refresh-token}, \texttt{logout} \\
Interno &
\texttt{/api/internals/validate-token} \\
Usuarios &
CRUD \texttt{/api/users} \\
Roles &
CRUD \texttt{/api/roles}; asignación de users/modules/menus/permissions \\
Módulos &
CRUD \texttt{/api/modules} \\
Menús &
CRUD \texttt{/api/menus}; \texttt{GET /api/menus/tree} \\
Permisos &
\texttt{GET /api/permissions} \\
\bottomrule
\end{tabular}
\end{center}

Nota de diseño: \texttt{GET /api/menus/tree} requiere sesión autenticada (rol en JWT)
y \textbf{no} el permiso administrativo \texttt{MENUS\_READ}, porque el árbol es
navegación por rol; \texttt{MENUS\_READ} queda para CRUD de menús.

% -----------------------------------------------------------------------------
\section{Frontend (SPA)}
% -----------------------------------------------------------------------------
La SPA implementa:
\begin{itemize}[leftmargin=1.5em]
  \item login y selector de rol con opción de cancelar / cambiar de usuario;
  \item almacenamiento de tokens en \texttt{sessionStorage};
  \item interceptor HTTP con Bearer y renovación de refresh;
  \item rutas protegidas y pantallas de 403, sesión expirada y token expirado;
  \item menú dinámico desde \texttt{/api/menus/tree};
  \item administración de usuarios y roles (creación, plantillas VENDEDOR/AUDITOR,
        asignación de permisos/módulos/menús/usuarios);
  \item pantalla de Ventas integrada.
\end{itemize}

La UI \textbf{no} toma decisiones finales de autorización: solo refleja permisos;
el backend siempre valida.

% -----------------------------------------------------------------------------
\section{Microservicio Ventas}
% -----------------------------------------------------------------------------
\begin{itemize}[leftmargin=1.5em]
  \item Sin base de datos de usuarios (almacén en memoria para demo académica).
  \item Endpoints: \texttt{GET /api/ventas} (\texttt{VENTAS\_READ}),
        \texttt{POST /api/ventas} (\texttt{VENTAS\_CREATE}).
  \item Middleware Zero Trust hacia el Master, con reintentos y backoff ante
        cold start típico de PaaS free tier (\texttt{503 MASTER\_UNAVAILABLE}).
\end{itemize}

Roles demo sembrados:
\begin{center}
\begin{tabular}{@{}lll@{}}
\toprule
\textbf{Usuario} & \textbf{Rol} & \textbf{Enfoque} \\
\midrule
\texttt{admin@example.com} & ADMIN (+ VENDEDOR, AUDITOR) & administración completa \\
\texttt{vendedor@example.com} & VENDEDOR & solo ventas \\
\texttt{auditor@example.com} & AUDITOR & lectura \\
\bottomrule
\end{tabular}
\end{center}
Password demo: \texttt{ChangeMe123!} (debe cambiarse en producción).

% -----------------------------------------------------------------------------
\section{Seguridad aplicada}
% -----------------------------------------------------------------------------
\begin{itemize}[leftmargin=1.5em]
  \item Hash de contraseñas con \textbf{Argon2}; nunca se expone \texttt{passwordHash}.
  \item Validación de fortaleza de password; mensajes de login genéricos
        (anti-enumeración).
  \item Helmet, CORS controlado, límites de body JSON.
  \item Rate limiting en login, refresh y validate-token.
  \item Validación de entrada con Zod; prevención de mass-assignment de
        \texttt{estado}.
  \item Soft delete en entidades principales.
  \item Secretos exclusivamente vía variables de entorno (\texttt{.env.example}).
  \item JWT con issuer/audience/expiración/jti; revocación en logout.
\end{itemize}

Decisión aceptada: firma simétrica HS256 + validación centralizada
(\texttt{validate-token}). El enunciado permite como alternativa RS256/ES256
en microservicios hijos; no es bloqueante para el cumplimiento funcional.

% -----------------------------------------------------------------------------
\section{DevSecOps y despliegue}
% -----------------------------------------------------------------------------
\subsection{Estrategia de ramas}
\begin{lstlisting}
feature/*  ->  dev  ->  test  ->  main
\end{lstlisting}
Política automatizada: \texttt{main} solo desde \texttt{test}; \texttt{test} solo desde \texttt{dev}.

\subsection{Pipeline (merge a main)}
\begin{enumerate}[leftmargin=1.5em]
  \item Instalación de dependencias, build y pruebas (Master, Frontend, Ventas).
  \item SonarCloud con Quality Gate obligatorio (\texttt{sonar.qualitygate.wait=true}).
  \item SAST adicional con modelo ML (CodeBERT-VulnCWE); hallazgos bloquean deploy.
  \item Despliegue vía Railway CLI.
  \item Notificaciones Telegram (inicio, pruebas, Sonar, ML, éxito/fallo, merges).
\end{enumerate}

\subsection{Despliegue recomendado}
Plataforma alineada al proyecto: \textbf{Railway} (alternativa: Render).
Servicios sugeridos: PostgreSQL + Master (+ Ventas + frontend estático).
Documentación operativa: \texttt{docs/DEVSECOPS.md}.

% -----------------------------------------------------------------------------
\section{Pruebas y verificación}
% -----------------------------------------------------------------------------
\begin{itemize}[leftmargin=1.5em]
  \item Pruebas unitarias con Vitest (suite del Master).
  \item Smoke E2E: \texttt{npm run smoke} (login, select-role, refresh, validate, logout, rechazo de token revocado).
  \item Verificación manual SPA: login $\rightarrow$ rol $\rightarrow$ Ventas $\rightarrow$ logout.
  \item Matriz de cumplimiento: \texttt{docs/COMPLIANCE\_CHECKLIST.md}.
\end{itemize}

Arranque local típico:
\begin{lstlisting}
npm run db:up
npm run db:reset
npm run dev
npm run ventas:dev
npm run frontend:dev
\end{lstlisting}
Nota: PostgreSQL en Docker usa puerto host \textbf{5433} para no chocar con un Postgres local en 5432.

% -----------------------------------------------------------------------------
\section{Estructura del repositorio}
% -----------------------------------------------------------------------------
\begin{lstlisting}
ProyectoP3_SWSeguro/
  src/                 # Master Gateway
  prisma/              # Schema, seed, seed demo roles
  frontend/            # SPA React
  services/ventas/     # Microservicio hijo
  .github/workflows/   # CI/CD y políticas de rama
  scripts/             # smoke + SAST ML
  docs/                # DevSecOps, checklist, informe
  README.md
\end{lstlisting}

% -----------------------------------------------------------------------------
\section{Resultados y conclusiones}
% -----------------------------------------------------------------------------
Se entregó un sistema Full-Stack seguro coherente con el enunciado académico:
Master Gateway como autoridad de identidad y autorización, SPA con selección
obligatoria de rol, microservicio Ventas bajo Zero Trust y pipeline DevSecOps
con controles Shift-Left.

\begin{itemize}[leftmargin=1.5em]
  \item \textbf{Cumplimiento funcional:} autenticación, RBAC, menú CTE, CRUD soft delete,
        Ventas, SPA y CI/CD.
  \item \textbf{Pendiente operativo (configuración externa):} secrets de GitHub,
        branch protection, proyecto SonarCloud y servicio Railway/Render.
  \item \textbf{Mejora opcional:} migración a JWT asimétrico (RS256/ES256) para
        validación offline en hijos.
\end{itemize}

% -----------------------------------------------------------------------------
\section{Anexos}
% -----------------------------------------------------------------------------
\subsection{A. Variables de entorno críticas (Master)}
\begin{itemize}[leftmargin=1.5em]
  \item \texttt{DATABASE\_URL}
  \item \texttt{JWT\_ACCESS\_SECRET}, \texttt{JWT\_TEMP\_SECRET}
  \item \texttt{INTERNAL\_API\_KEY}
  \item \texttt{JWT\_ISSUER}, \texttt{JWT\_AUDIENCE}
  \item TTL de access/temp/refresh
\end{itemize}

\subsection{B. Secrets de GitHub Actions (producción)}
\begin{itemize}[leftmargin=1.5em]
  \item \texttt{TELEGRAM\_BOT\_TOKEN}, \texttt{TELEGRAM\_CHAT\_ID}
  \item \texttt{SONAR\_TOKEN}, \texttt{SONAR\_PROJECT\_KEY}, \texttt{SONAR\_ORGANIZATION}
  \item \texttt{RAILWAY\_TOKEN}, \texttt{RAILWAY\_SERVICE}, \texttt{APP\_URL}
\end{itemize}

\subsection{C. Cómo incorporar este informe a otro documento \LaTeX}
\begin{enumerate}[leftmargin=1.5em]
  \item Opción simple: copiar desde \verb|\section{Introducción}| hasta el final
        (sin \verb|\documentclass| ni \verb|\begin{document}|).
  \item Opción modular: en el documento padre usar
        \begin{lstlisting}
% Cuerpo del informe (sin preámbulo) para incorporar en otro documento LaTeX:
%   % Cuerpo del informe (sin preámbulo) para incorporar en otro documento LaTeX:
%   % Cuerpo del informe (sin preámbulo) para incorporar en otro documento LaTeX:
%   \input{docs/informe_proyecto_cuerpo}
% Requiere en el documento padre (según uso): babel spanish, booktabs, listings,
% enumitem, hyperref (opcionales pero recomendados).

\section{Introducción}
El proyecto académico \emph{Sistema de Autenticación y Autorización Centralizado
(Master Gateway)} responde a la necesidad de una puerta maestra de seguridad
para arquitecturas de microservicios. En lugar de que cada servicio gestione
usuarios y credenciales, el Master autentica usuarios, obliga a seleccionar un
rol activo, emite JWT con privilegio mínimo, expone validación interna Zero Trust
para microservicios hijos y administra usuarios, roles, módulos, menús y permisos
con soft delete y auditoría.

Principios rectores: Zero Trust, Least Privilege, Shift-Left Security,
autenticación stateless y soft delete.

\section{Objetivos}
\begin{enumerate}
  \item Implementar el flujo login $\rightarrow$ TempToken $\rightarrow$ selección de rol
        $\rightarrow$ AccessToken + RefreshToken rotativo.
  \item Modelar entidades con auditoría y estado ACTIVO/INACTIVO.
  \item Generar menú dinámico con CTE recursiva filtrado por rol JWT.
  \item Integrar un microservicio hijo (Ventas) bajo Zero Trust.
  \item Entregar SPA con rutas protegidas y navegación dinámica.
  \item Automatizar calidad y despliegue (GitHub Actions, SonarCloud, SAST/ML, Railway, Telegram).
\end{enumerate}

\section{Arquitectura del sistema}
\subsection{Vista lógica}
\begin{enumerate}
  \item \textbf{Master Gateway}: API Node.js/TypeScript (auth, RBAC, CRUD, validate-token).
  \item \textbf{PostgreSQL} vía Prisma ORM.
  \item \textbf{SPA React + Vite}: login, selector de rol, administración y Ventas.
  \item \textbf{Microservicio Ventas}: sin DB de usuarios; valida tokens en el Master.
\end{enumerate}

\subsection{Flujo Zero Trust}
El frontend envía el AccessToken a Ventas; Ventas llama a
\texttt{POST /api/internals/validate-token} con \texttt{x-internal-api-key}.
Solo si el Master confirma validez y permisos se ejecuta la lógica de negocio.

\subsection{Stack}
Master: Express, TypeScript, Prisma, Argon2, Jose (JWT), Helmet, Zod.
Frontend: React 19, Vite, React Router.
CI/CD: GitHub Actions, SonarCloud, SAST/ML, Railway CLI, Telegram.

\section{Modelo de datos}
Entidades: User, Role, UserRole, Module, RoleModule, Menu (Adjacency List),
RoleMenu, Permission, RolePermission, RefreshToken, TokenRevocation, AuditLog.
Menú: solo hojas con URL, anticiclos, exclusión de inactivos.

\section{Autenticación y autorización}
Login emite TempToken + roles; la SPA fuerza select-role; el AccessToken incluye
únicamente el rol seleccionado y sus permisos. Refresh rotativo y logout con
revocación por \texttt{jti}. Endpoint interno
\texttt{/api/internals/validate-token} responde
\texttt{active}, \texttt{userId}, \texttt{roleId}, \texttt{roleName}, \texttt{permissions}.

\section{Frontend}
SPA con sessionStorage, interceptor Bearer/refresh, menú dinámico, administración
de usuarios/roles, pantallas 403/sesión/token expirado y cancelación del selector
de rol para cambiar de usuario. La UI no decide autorización final.

\section{Microservicio Ventas}
Endpoints con permisos \texttt{VENTAS\_READ}/\texttt{VENTAS\_CREATE}; retries ante
cold start PaaS. Roles demo: ADMIN, VENDEDOR, AUDITOR.

\section{Seguridad}
Argon2, rate limiting, Helmet, Zod, soft delete, secretos en entorno, mensajes
de login genéricos. JWT HS256 + validate-token (RS256 opcional no bloqueante).

\section{DevSecOps}
Ramas \texttt{feature $\rightarrow$ dev $\rightarrow$ test $\rightarrow$ main}.
Pipeline en merge a main: build/tests, Sonar Quality Gate obligatorio, SAST/ML,
deploy Railway, Telegram. Despliegue recomendado: Railway (alternativa Render).

\section{Pruebas}
Vitest, smoke E2E (\texttt{npm run smoke}), checklist en
\texttt{docs/COMPLIANCE\_CHECKLIST.md}.

\section{Conclusiones}
Se cumple el alcance académico Full-Stack seguro (Master + SPA + Ventas + DevSecOps).
Pendiente operativo: secrets GitHub, branch protection, SonarCloud y PaaS.
Mejora opcional: JWT asimétrico.

% Requiere en el documento padre (según uso): babel spanish, booktabs, listings,
% enumitem, hyperref (opcionales pero recomendados).

\section{Introducción}
El proyecto académico \emph{Sistema de Autenticación y Autorización Centralizado
(Master Gateway)} responde a la necesidad de una puerta maestra de seguridad
para arquitecturas de microservicios. En lugar de que cada servicio gestione
usuarios y credenciales, el Master autentica usuarios, obliga a seleccionar un
rol activo, emite JWT con privilegio mínimo, expone validación interna Zero Trust
para microservicios hijos y administra usuarios, roles, módulos, menús y permisos
con soft delete y auditoría.

Principios rectores: Zero Trust, Least Privilege, Shift-Left Security,
autenticación stateless y soft delete.

\section{Objetivos}
\begin{enumerate}
  \item Implementar el flujo login $\rightarrow$ TempToken $\rightarrow$ selección de rol
        $\rightarrow$ AccessToken + RefreshToken rotativo.
  \item Modelar entidades con auditoría y estado ACTIVO/INACTIVO.
  \item Generar menú dinámico con CTE recursiva filtrado por rol JWT.
  \item Integrar un microservicio hijo (Ventas) bajo Zero Trust.
  \item Entregar SPA con rutas protegidas y navegación dinámica.
  \item Automatizar calidad y despliegue (GitHub Actions, SonarCloud, SAST/ML, Railway, Telegram).
\end{enumerate}

\section{Arquitectura del sistema}
\subsection{Vista lógica}
\begin{enumerate}
  \item \textbf{Master Gateway}: API Node.js/TypeScript (auth, RBAC, CRUD, validate-token).
  \item \textbf{PostgreSQL} vía Prisma ORM.
  \item \textbf{SPA React + Vite}: login, selector de rol, administración y Ventas.
  \item \textbf{Microservicio Ventas}: sin DB de usuarios; valida tokens en el Master.
\end{enumerate}

\subsection{Flujo Zero Trust}
El frontend envía el AccessToken a Ventas; Ventas llama a
\texttt{POST /api/internals/validate-token} con \texttt{x-internal-api-key}.
Solo si el Master confirma validez y permisos se ejecuta la lógica de negocio.

\subsection{Stack}
Master: Express, TypeScript, Prisma, Argon2, Jose (JWT), Helmet, Zod.
Frontend: React 19, Vite, React Router.
CI/CD: GitHub Actions, SonarCloud, SAST/ML, Railway CLI, Telegram.

\section{Modelo de datos}
Entidades: User, Role, UserRole, Module, RoleModule, Menu (Adjacency List),
RoleMenu, Permission, RolePermission, RefreshToken, TokenRevocation, AuditLog.
Menú: solo hojas con URL, anticiclos, exclusión de inactivos.

\section{Autenticación y autorización}
Login emite TempToken + roles; la SPA fuerza select-role; el AccessToken incluye
únicamente el rol seleccionado y sus permisos. Refresh rotativo y logout con
revocación por \texttt{jti}. Endpoint interno
\texttt{/api/internals/validate-token} responde
\texttt{active}, \texttt{userId}, \texttt{roleId}, \texttt{roleName}, \texttt{permissions}.

\section{Frontend}
SPA con sessionStorage, interceptor Bearer/refresh, menú dinámico, administración
de usuarios/roles, pantallas 403/sesión/token expirado y cancelación del selector
de rol para cambiar de usuario. La UI no decide autorización final.

\section{Microservicio Ventas}
Endpoints con permisos \texttt{VENTAS\_READ}/\texttt{VENTAS\_CREATE}; retries ante
cold start PaaS. Roles demo: ADMIN, VENDEDOR, AUDITOR.

\section{Seguridad}
Argon2, rate limiting, Helmet, Zod, soft delete, secretos en entorno, mensajes
de login genéricos. JWT HS256 + validate-token (RS256 opcional no bloqueante).

\section{DevSecOps}
Ramas \texttt{feature $\rightarrow$ dev $\rightarrow$ test $\rightarrow$ main}.
Pipeline en merge a main: build/tests, Sonar Quality Gate obligatorio, SAST/ML,
deploy Railway, Telegram. Despliegue recomendado: Railway (alternativa Render).

\section{Pruebas}
Vitest, smoke E2E (\texttt{npm run smoke}), checklist en
\texttt{docs/COMPLIANCE\_CHECKLIST.md}.

\section{Conclusiones}
Se cumple el alcance académico Full-Stack seguro (Master + SPA + Ventas + DevSecOps).
Pendiente operativo: secrets GitHub, branch protection, SonarCloud y PaaS.
Mejora opcional: JWT asimétrico.

% Requiere en el documento padre (según uso): babel spanish, booktabs, listings,
% enumitem, hyperref (opcionales pero recomendados).

\section{Introducción}
El proyecto académico \emph{Sistema de Autenticación y Autorización Centralizado
(Master Gateway)} responde a la necesidad de una puerta maestra de seguridad
para arquitecturas de microservicios. En lugar de que cada servicio gestione
usuarios y credenciales, el Master autentica usuarios, obliga a seleccionar un
rol activo, emite JWT con privilegio mínimo, expone validación interna Zero Trust
para microservicios hijos y administra usuarios, roles, módulos, menús y permisos
con soft delete y auditoría.

Principios rectores: Zero Trust, Least Privilege, Shift-Left Security,
autenticación stateless y soft delete.

\section{Objetivos}
\begin{enumerate}
  \item Implementar el flujo login $\rightarrow$ TempToken $\rightarrow$ selección de rol
        $\rightarrow$ AccessToken + RefreshToken rotativo.
  \item Modelar entidades con auditoría y estado ACTIVO/INACTIVO.
  \item Generar menú dinámico con CTE recursiva filtrado por rol JWT.
  \item Integrar un microservicio hijo (Ventas) bajo Zero Trust.
  \item Entregar SPA con rutas protegidas y navegación dinámica.
  \item Automatizar calidad y despliegue (GitHub Actions, SonarCloud, SAST/ML, Railway, Telegram).
\end{enumerate}

\section{Arquitectura del sistema}
\subsection{Vista lógica}
\begin{enumerate}
  \item \textbf{Master Gateway}: API Node.js/TypeScript (auth, RBAC, CRUD, validate-token).
  \item \textbf{PostgreSQL} vía Prisma ORM.
  \item \textbf{SPA React + Vite}: login, selector de rol, administración y Ventas.
  \item \textbf{Microservicio Ventas}: sin DB de usuarios; valida tokens en el Master.
\end{enumerate}

\subsection{Flujo Zero Trust}
El frontend envía el AccessToken a Ventas; Ventas llama a
\texttt{POST /api/internals/validate-token} con \texttt{x-internal-api-key}.
Solo si el Master confirma validez y permisos se ejecuta la lógica de negocio.

\subsection{Stack}
Master: Express, TypeScript, Prisma, Argon2, Jose (JWT), Helmet, Zod.
Frontend: React 19, Vite, React Router.
CI/CD: GitHub Actions, SonarCloud, SAST/ML, Railway CLI, Telegram.

\section{Modelo de datos}
Entidades: User, Role, UserRole, Module, RoleModule, Menu (Adjacency List),
RoleMenu, Permission, RolePermission, RefreshToken, TokenRevocation, AuditLog.
Menú: solo hojas con URL, anticiclos, exclusión de inactivos.

\section{Autenticación y autorización}
Login emite TempToken + roles; la SPA fuerza select-role; el AccessToken incluye
únicamente el rol seleccionado y sus permisos. Refresh rotativo y logout con
revocación por \texttt{jti}. Endpoint interno
\texttt{/api/internals/validate-token} responde
\texttt{active}, \texttt{userId}, \texttt{roleId}, \texttt{roleName}, \texttt{permissions}.

\section{Frontend}
SPA con sessionStorage, interceptor Bearer/refresh, menú dinámico, administración
de usuarios/roles, pantallas 403/sesión/token expirado y cancelación del selector
de rol para cambiar de usuario. La UI no decide autorización final.

\section{Microservicio Ventas}
Endpoints con permisos \texttt{VENTAS\_READ}/\texttt{VENTAS\_CREATE}; retries ante
cold start PaaS. Roles demo: ADMIN, VENDEDOR, AUDITOR.

\section{Seguridad}
Argon2, rate limiting, Helmet, Zod, soft delete, secretos en entorno, mensajes
de login genéricos. JWT HS256 + validate-token (RS256 opcional no bloqueante).

\section{DevSecOps}
Ramas \texttt{feature $\rightarrow$ dev $\rightarrow$ test $\rightarrow$ main}.
Pipeline en merge a main: build/tests, Sonar Quality Gate obligatorio, SAST/ML,
deploy Railway, Telegram. Despliegue recomendado: Railway (alternativa Render).

\section{Pruebas}
Vitest, smoke E2E (\texttt{npm run smoke}), checklist en
\texttt{docs/COMPLIANCE\_CHECKLIST.md}.

\section{Conclusiones}
Se cumple el alcance académico Full-Stack seguro (Master + SPA + Ventas + DevSecOps).
Pendiente operativo: secrets GitHub, branch protection, SonarCloud y PaaS.
Mejora opcional: JWT asimétrico.

        \end{lstlisting}
        tras mover el cuerpo a un archivo sin preámbulo.
  \item Para Word/Google Docs: compilar a PDF y convertir, o usar Pandoc:
        \begin{lstlisting}
pandoc docs/informe_proyecto.tex -o informe.docx
        \end{lstlisting}
\end{enumerate}

\end{document}
